\documentclass[12pt,a4paper]{ctexart}

% ========== 页面设置 ==========
\usepackage[a4paper,margin=2.5cm]{geometry}
\usepackage{setspace}
\usepackage{enumitem}
\usepackage{graphicx}  % 插入图片的宏包
\usepackage{titlesec}
\usepackage{amsmath,amssymb}
\usepackage{fontspec}
\usepackage{fancyhdr}          % 控制页眉页脚
\usepackage[hidelinks]{hyperref}
\usepackage[backend=biber,style=ieee]{biblatex}

% bib
\addbibresource{references.bib}

% ========== 字体设置 ==========
\setmainfont{Times New Roman} % 英文主字体
\setCJKmainfont{SimSun}       % 中文主字体：宋体

% ========== 段落与格式 ==========
\setlength{\parindent}{2em}
\setlength{\parskip}{0.8em}
\linespread{1.3}  % 行距

% ========== 标题格式 ==========
\titleformat{\section}
  {\bfseries\large}
  {\thesection、}
  {0.5em}
  {}

\titleformat{\subsection}
  {\normalsize}
  {\thesubsection\quad}
  {0em}
  {}

% 页眉，页脚
\pagestyle{fancy}                   % 启用 fancy 样式
\fancyhf{}                          % 清空原有页眉页脚
\fancyfoot[C]{\thepage}             % 页脚中间显示页码
\renewcommand{\headrulewidth}{0pt}  % 去掉页眉横线
\renewcommand{\footrulewidth}{0pt}  % 去掉页脚横线（可选）

% ========== 正文 ==========
\begin{document}

\begin{center}
    \LARGE \textbf{高等数字集成电路作业-第三章}
\end{center}

\section{基础概念问题}

\subsection{请简要描述基于VerilogHDL语言对VLSI/FPGA设计带来的意义?}

VerilogHDL提供了行为级的数字逻辑电路描述方式，使设计者能在高层抽象下描述功能，便于数字电路的模块化。支持模块设计对于大型的VLSI设计尤其重要，极大的提高的设计效率。对于FPGA设计，相比于ASIC平台，使用FPGA实现数字电路是一个更轻便的方案，需要高效的硬件描述和快速的原型设计。VerilogHDL以基于c语法的设计思想，提供了一种简洁的方式来描述FPGA中的数字电路，并可以通过综合工具将设计直接映射到FPGA上。

\subsection{请简要描述基于VerilogHDL语言完成的可综合电路与不可综合电路的特点？}

可综合与不可综合的本质是，使用Verilog描述的结构，能否映射到可实现的电路。可综合的实际电路由组合逻辑和时序逻辑组成，通过Verilog中的assign与always语句块进行描述。不可综合电路最大的特点，就是误将用于仿真的语法结构应用在电路描述中，如延迟赋值、系统任务函数、赋值X态等，退出条件不定的循环等。这些情况都没有实际的意义，比如在verilog可以赋值x态，在仿真中可以明显看到x态，但是在实际的电路中，寄存器处于x态是亚稳态，没有意义。

\subsection{查阅文献，介绍CMOS触发器的电路发明简史，分析CMOS触发器电路引入对构建基于标准单元的超大规模数字集成电路设计的底层电路的意义和作用？描述非阻塞verilog语句是如何对触发器进行建模的？分析用非阻塞verilog语句进行触发器建模的优势？}

最早的触发器概念源于早期电子开关电路。1918年William Eccles与F.W. Jordan发明了Eccles‑Jordan触发电路作为一种二态存储元件\cite{eccles}，使用的是双稳态继电器或电子管的储存结构。随着晶体管和MOSFET的发展，触发器的设计只是在前者的基础上，改为MOSFET作为存储元件，实现MOS技术下的锁存器、触发器。但随着1963年，Frank Wanlass 与 Chih-Tang Sah在1963年的ISSCC会议上阐述CMOS方案\cite{cmos}，触发器的结构有了新的实现方案，触发器及锁存器结构得以采用传输门实现。

触发器作为同步数字系统中关键的状态存储元件，使得寄存器‐传输级设计能够通过标准单元库进行模块化、可复用的实现。CMOS技术本身具备低静态功耗、高密度集成等特点，这使得将触发器设计成标准单元库中的一个单元变得可行，从而加速了从RTL到综合的过程。触发器是一切时序逻辑的基础，是VLSI的底层支柱。

Verilog用非阻塞语句表示时序逻辑，对应了触发器的物理表现。其将触发器两次采样间的时钟周期窗口建模为赋值开始到赋值结束，在窗口结束之前触发器的输入信号是允许改变的，但最终触发器的采样结果只由赋值结束时（触发器的有效采样窗口）的数据决定。

使用非阻塞语句对触发器建模，真实地反映硬件中多个触发器在同一时钟边沿“同时”触发、更新状态的并行特性，有利于仿真、综合等软件程序对底层硬件的理解。

\subsection{请简要描述阻塞描述语句与非阻塞描述语句的各自适用电路与使用注意事项？}

阻塞赋值语句中，赋值操作的右值会在当前仿真时间步立即计算，并赋给信号。在实际电路中，用于描述组合逻辑电路，因为这些电路没有时序依赖，只是对信号逻辑路径的一种描述。在使用时，需要注意避免组合逻辑输出状态不全的情况下，在该情况下综合器会认为组合逻辑电路带有记忆功能，但不使用触发器，而是锁存器实现，无意识的使用锁存器会导致严重的问题。同时，组合逻辑电路要避免环回驱动，及组合逻辑中，信号路径不能出现环回，输入不能来自输出。

非阻塞赋值语句中，会在当前仿真时间步结束时才将结果更新给信号。在实际电路中，用于描述时序逻辑电路。在使用中尤其要注意，对同一个触发器，不能有多个电路同时进行赋值，也就是Multi-drive问题。

\subsection{请简要分析全同步电路与异步电路在常规VLSI电路设计过程中的设计应用注意事项与考虑思路？}

全同步电路依赖于全局时钟信号，所有的操作都由时钟信号同步控制，这使得时序控制变得简单，设计和调试更为直观。然而，时钟延迟和时钟偏移问题可能影响电路性能，需要特别关注时钟分配和时序分析。此外，高频时钟会带来功耗问题，功耗优化问题尤为重要。相比之下，异步电路没有全局统一时钟，它通过数据之间的相互作用进行触发操作，设计上具有更大的灵活性。异步电路适合低功耗、低时延的应用，但设计与验证较为复杂，异步电路间的同步存在不稳定的问题，通常使用握手信号和延迟锁存器来保证数据传输的稳定性，但由于没有全局时钟，数据同步和时序验证依然复杂。选择同步或异步电路取决于具体应用对功耗、速度、设计复杂度以及验证的要求。

\subsection{请简要分析Latch与D-Flip-Flop电路在常规VLSI电路设计过程中的设计应用注意事项与考虑思路？}

锁存器与触发器同样是存储元件，由于其只有一级传输门，构造简单，传输延迟更小，相对的是在时钟有效时对输入信号进行更新（锁存），易受竞态条件影响。面积小，功耗相对较低，在高频或异步设计中会采用锁存器。而触发器一般由两级传输门组成，是基于时钟同步设计的基础单元，在时序逻辑电路中广泛应用，具有较强的抗干扰能力和时序可控性，适用于要求高可靠性和同步操作的设计，虽然面积和功耗相比触发器更大，但电路设计更加可靠、简单。

\printbibliography

\end{document}

